\section{Hilbert series and fixed points of the grading}\label{sec:hilbert}

Throughout, $R=\bigoplus_{d\ge0}R_d$ is a finitely generated positively graded
$\C$-algebra with $R_0=\C$, and $\Hilb_R(t)=\sum_d\dim R_d\,t^d$. It is a
rational function whose poles are roots of unity. For a root of unity $\eta$,
the grading element $\eta$ acts on $R_d$ by $\eta^d$; we write $X=\Spec R$ and
$X^\eta$ for the scheme-theoretic fixed locus. If $\eta$ has order $q$, then
$X^\eta=\Spec R/I_q$ with
\begin{equation}\label{eq:fixedideal}
 I_q=(R_d:\ q\nmid d),
\end{equation}
since a homogeneous element of degree $d$ is changed by the factor $\eta^d$.

\begin{lemma}\label{lem:isolated}
If $X^\eta$ is supported at the graded origin, then $\Hilb_R(t)$ is regular at
every primitive root of unity of the order $q$ of $\eta$.
\end{lemma}

\begin{proof}
Choose homogeneous generators $a_1,\dots,a_l$ of $I_q$ with degrees $e_i$ not
divisible by $q$. As the support is the origin, $R/I_q$ is finite-dimensional.
The homology modules $N_j$ of the graded Koszul complex $K_R(a_1,\dots,a_l)$
are finitely generated over $R$ and annihilated by $I_q$, because
multiplication by each $a_i$ is null-homotopic on the Koszul complex; hence
they are finite-dimensional. The graded Euler characteristic gives
\[
 \Hilb_R(t)\prod_{i=1}^l(1-t^{e_i})=\sum_j(-1)^j\Hilb_{N_j}(t),
\]
whose right side is a polynomial. No factor $1-t^{e_i}$ vanishes at a
primitive $q$-th root of unity.
\end{proof}

\begin{proposition}\label{prop:fixedcomponentpole}
Let $R$ be a domain, $M\ge2$, $\zeta$ a primitive $M$-th root of unity. Suppose
$X^\zeta$ has a unique irreducible component $Y$ of maximal dimension
$\delta\ge1$, and $X$ is smooth at the generic point of $Y$. Give $Y$ its
reduced structure; its coordinate ring is concentrated in degrees divisible
by $M$; put $h_Y(z)=\sum_k\dim\C[Y]_{Mk}z^k$ and
$e_Y=\lim_{z\to1}(1-z)^\delta h_Y(z)>0$. Then the conormal space of $Y$ in $X$
at its generic point has no eigenvalue one under $\zeta$, and if
$\nu_1,\dots,\nu_c$ ($c=\dim X-\delta$) are its eigenvalues,
\begin{equation}\label{eq:residue}
 \lim_{t\to\zeta}(1-t^M)^\delta\Hilb_R(t)=\frac{e_Y}{\prod_{j=1}^c(1-\nu_j)}\ne0 .
\end{equation}
In particular $\Hilb_R$ has a pole of order exactly $\delta$ at $\zeta$, and
every expression $\Hilb_R(t)=P(t)/\prod_i(1-t^{a_i})^{b_i}$ with $P\in\C[t]$
has some $a_i$ divisible by $M$.
\end{proposition}

\begin{proof}
Choose homogeneous algebra generators of $R$ and list those of degree not
divisible by $M$ as $f_1,\dots,f_l$, $\deg f_i=w_i$. Their ideal $I$ is the
fixed-locus ideal: $f_i$ is a nonzero multiple of $\zeta\cdot f_i-f_i$, and
modulo the $f_i$ every remaining generator is fixed. Thus $X^\zeta=\Spec R/I$
and $(R/I)_d=0$ when $M\nmid d$. The ideal is homogeneous, so every component
of the fixed locus is stable under the connected grading torus, hence graded;
the coordinate ring of $Y$ lives in degrees divisible by $M$.

Let $K_\bullet$ be the graded Koszul complex on $f_1,\dots,f_l$, its $i$-th
exterior generator of degree $w_i$ and $\zeta$-character $\zeta^{w_i}$, and
$N_j$ its homology. As before, each $N_j$ is a finite graded $R/I$-module, and
\begin{equation}\label{eq:euler}
 \Hilb_R(t)\prod_{i=1}^l(1-t^{w_i})=\sum_{j=0}^l(-1)^j\Hilb_{N_j}(t).
\end{equation}
All factors $1-t^{w_i}$ are nonzero at $\zeta$.

Every finite graded $R/I$-module $N$ splits into its degree classes modulo
$M$; each class, regraded so that degree $M$ becomes degree one, has support
of dimension at most $\delta$, so $\Hilb_N$ has a pole of order at most
$\delta$ at $\zeta$. If the localization $N_\eta$ at the generic point $\eta$
of $Y$ is a $\C(Y)$-vector space, then
\begin{equation}\label{eq:leading}
 \lim_{t\to\zeta}(1-t^M)^\delta\Hilb_N(t)=e_Y\Tr(\zeta\mid N_\eta).
\end{equation}
Indeed, take a graded prime filtration of each degree class of $N$. Factors
supported in dimension less than $\delta$ contribute zero to the limit. Since
$Y$ is the unique component of dimension $\delta$, the remaining factors are
shifted copies of $\C[Y]$; one generated in degree $a$ contributes
$t^ah_Y(t^M)$, with leading coefficient $\zeta^ae_Y$. After localization
these factors give a filtration of $N_\eta$ by lines of character $\zeta^a$,
which gives the trace. That $e_Y>0$ is the positivity of the leading Hilbert
coefficient of the positively graded domain $\C[Y]$ regraded by $M$.

At a smooth closed point of $X$ on $Y$, diagonalize the finite-order action of
$\zeta$ on the cotangent space and lift eigenvectors to the maximal ideal by
the averaging projectors. They are formal coordinates, the completed local
ring is the power series ring in them with diagonal action, and the fixed
scheme is cut out by the coordinates with eigenvalue different from one. So
the fixed scheme is smooth and regularly embedded there, and its conormal
representation has no eigenvalue one. This applies at the generic point of
$Y$.

Put $K=\C(Y)$. In the regular local ring $R_\eta$, the ideal $I_\eta$ is the
maximal ideal, generated by a regular sequence of length $c$. There is an
exact sequence of $K$-representations of $\langle\zeta\rangle$
\[
 0\to L\to E=\bigoplus_{i=1}^lKe_i\to V=I_\eta/I_\eta^2\to0,\qquad
 e_i\mapsto f_i\bmod I_\eta^2 ,
\]
where $E$ has the characters $\zeta^{w_i}\ne1$ and $V$ the characters
$\nu_1,\dots,\nu_c\ne1$. Choosing among the $f_i$ a minimal generating set of
$I_\eta$, which is a regular sequence, and subtracting combinations of it from
the other generators changes the Koszul complex into the Koszul complex of a
regular sequence tensored with the Koszul complex on $l-c$ zeros; hence
$(N_j)_\eta\cong\bigwedge^j_KL$, naturally and therefore equivariantly. These
modules are vector spaces, so \eqref{eq:leading} applies, and
$\sum_j(-1)^j\Tr(\zeta\mid(N_j)_\eta)=\det(1-\zeta\mid L)$. Multiplying
\eqref{eq:euler} by $(1-t^M)^\delta$ and letting $t\to\zeta$,
\[
 \lim_{t\to\zeta}(1-t^M)^\delta\Hilb_R(t)
 =e_Y\frac{\det(1-\zeta\mid L)}{\prod_i(1-\zeta^{w_i})}
 =\frac{e_Y}{\det(1-\zeta\mid V)} ,
\]
which is \eqref{eq:residue}. The pole order is at most $\delta$ by the module
argument and exactly $\delta$ since the limit is not zero. Finally, if no
$a_i$ were divisible by $M$, the expression $P/\prod_i(1-t^{a_i})^{b_i}$
would be regular at $\zeta$.
\end{proof}

The Koszul lemmas used (null-homotopy, change of generators, vanishing of
higher Koszul homology for a regular sequence) are
\cite[Tags 0621, 062F]{Stacks}. Nonreduced structure of $X^\zeta$ at the
origin does no harm: its contributions are finite-dimensional.

\begin{lemma}\label{lem:orbitclosure}
Let $p\in X$ be a closed point other than the origin, and suppose that its
stabilizer in the grading group $\Gm$ is exactly $\mu_M$. Let $Y$ be the
closure of its $\Gm$-orbit, with reduced structure. Then $\C[Y]$ is
concentrated in degrees divisible by $M$, the set $\{k:\C[Y]_{Mk}\ne0\}$ is a
numerical semigroup (cofinite in $\Z_{\ge0}$), each nonzero
$\C[Y]_{Mk}$ is a line, and $e_Y=1$ with $\delta=1$.
\end{lemma}

\begin{proof}
Restriction to the orbit is the homomorphism $R\to\C[u]$, $f\mapsto f(p)u^d$
for $f\in R_d$; its image is $\C[Y]$, spanned by the $u^d$ with $d\in S=\{d:
f(p)\ne0\text{ for some }f\in R_d\}$, a semigroup since products of functions
nonzero at $p$ are nonzero at $p$. An element $\lambda\in\Gm$ fixes $p$ if and
only if $f(\lambda\cdot p)=\lambda^df(p)=f(p)$ for all homogeneous $f$, since
homogeneous functions separate the points of $X$; that is, if and only if
$\lambda^d=1$ for all $d\in S$. So the stabilizer is $\mu_{\gcd S}$, and
$\gcd S=M$. A subsemigroup of $M\Z_{\ge0}$ whose elements have greatest
common divisor $M$ contains all sufficiently large multiples of $M$. Hence
$h_Y(z)=1/(1-z)-P(z)$ for a polynomial $P$, and $e_Y=1$.
\end{proof}

\begin{lemma}\label{lem:lift}
Let $GL_n$ act linearly on a vector space $U$, let $R=\C[U]^{SL_n}$ be graded
so that $R_0=\C$, and let $\gamma$ be an odd integer such that every
$B\in GL_n$ acts on $X=\Spec R$ as the grading element $(\det B)^\gamma$. Let
$p\in X$ be a closed point other than the origin, fixed by a grading element
$\eta$ of order $2^a$, and let $y\in U$ lie in the closed $SL_n$-orbit over
$p$. Then there is a diagonalizable $A\in GL_n$ of finite two-power order with
$A\cdot y=y$ and $\ord(\det A)=2^a$.
\end{lemma}

\begin{proof}
As $\gamma$ is odd, there is $\theta$ of order $2^a$ with $\theta^\gamma=\eta$.
Choose $B\in GL_n$ with $\det B=\theta$. Then $B\cdot y$ lies in a closed orbit
over $\eta\cdot p=p$, and each fibre of the quotient map contains a unique
closed orbit \cite[Section I.1, p.~84]{Luna}; so $B\cdot y=S\cdot y$ with
$S\in SL_n$, and $C=S^{-1}B$ fixes $y$ with $\det C=\theta$. The semisimple
part $C_s$ of $C$ also fixes $y$ (an eigenvector of $C_s$ for an eigenvalue
$\ne1$ cannot be fixed by $C=C_sC_u$, whose unipotent part preserves the
eigenspaces of $C_s$) and has the same determinant. Diagonalize $C_s$ and let
$\Gamma\subset\C^*$ be the group generated by its eigenvalues, a finitely
generated abelian group, the product of its finite cyclic torsion subgroup and
a free group. The projection of $\Gamma$ onto the two-primary part of its
torsion subgroup, trivial on the free factor, is a homomorphism
$\Gamma\to\C^*$; applied to every eigenvalue it preserves every multiplicative
relation among them, in particular the relations expressing $C_s\cdot y=y$
weight by weight, and it fixes $\theta$. The resulting diagonal matrix $A$
fixes $y$, has finite two-power order, and $\det A=\theta$.
\end{proof}

\begin{lemma}\label{lem:cyclotomic}
Let $p$ be a prime and $d=p^re$ with $p\nmid e$. In $\F_p[t]$,
$\overline{\Phi_d}=\overline{\Phi_e}$ if $r=0$ and
$\overline{\Phi_d}=\overline{\Phi_e}^{\,p^{r-1}(p-1)}$ if $r\ge1$. The roots of
$\overline{\Phi_e}$ in $\overline\F_p$ are the elements of multiplicative order
$e$, each simple.
\end{lemma}

\begin{proof}
For $r\ge1$, $\Phi_{ep^r}(t)\Phi_e(t^{p^{r-1}})=\Phi_e(t^{p^r})$ in $\Z[t]$
(group the roots on the right by the $p$-part of their order). Modulo $p$,
$t\mapsto t^{p^j}$ becomes the $p^j$-th power; cancel
$\overline{\Phi_e}^{\,p^{r-1}}$. The factorization $t^e-1=\prod_{b\mid e}
\Phi_b(t)$ persists modulo $p$, $t^e-1$ is separable there, and induction on
$e$ assigns to $\overline{\Phi_e}$ exactly the roots of order $e$.
\end{proof}

\begin{theorem}\label{thm:transfer}
Let $F,G\in\Q(t)\cap\Z[[t]]$ have all their poles at roots of unity, let $p$ be
a prime, $p\nmid M$, and suppose $F\equiv G\pmod{p\Z[[t]]}$. Suppose $G$ has a
pole of order $\delta\ge1$ at a primitive $M$-th root of unity $\zeta$, no pole
at any primitive $Mp^r$-th root of unity for $r\ge1$, and that
$C=\lim_{t\to\zeta}(1-t^M)^\delta G(t)$ is a unit in $\Z[\zeta]_{\mathfrak p}$
for some prime $\mathfrak p$ over $p$. Then $F$ has a pole at a primitive
$Mp^r$-th root of unity for some $r\ge0$, hence at all of them; and every
expression $F(t)=P(t)/\prod_i(1-t^{a_i})^{b_i}$ with $P\in\Q[t]$ and $a_i,b_i>0$
has some $a_i$ divisible by $Mp^r$, in particular $a_i\ge M$.
\end{theorem}

\begin{proof}
A reduced expression of $F$ has a denominator that is, up to sign, a product
of cyclotomic polynomials; normalize its constant term to $1$. The numerator is
then integral, being the product of the integral series $F$ with an integral
polynomial. So $F=N_F/D_F$ and $G=N_G/D_G$ with $N_\bullet,D_\bullet\in\Z[t]$,
$D_\bullet(0)=1$, coprime over $\Q$, the $D_\bullet$ signed products of
cyclotomic polynomials. Write $D_G=\varepsilon\Phi_M^\delta D_0$. By
Lemma~\ref{lem:cyclotomic} and the hypothesis on poles of order $Mp^r$, the
reduction $\overline{D_0}$ does not vanish at $\alpha=\zeta\bmod\mathfrak p$,
which has exact order $M$. With $U_M=(1-t^M)/\Phi_M$, separability of
$t^M-1$ modulo $p$ gives $\overline{U_M}(\alpha)\ne0$. So $D_0(\zeta)$ and
$U_M(\zeta)$ are $\mathfrak p$-units, and
$C=U_M(\zeta)^\delta N_G(\zeta)/(\varepsilon D_0(\zeta))$ forces
$N_G(\zeta)$ to be a unit: $\overline{N_G}(\alpha)\ne0$ while
$\overline{D_G}$ vanishes at $\alpha$. So $\overline G\in\F_p(t)$ has a pole
at $\alpha$. Reduction of both fractions is defined in $\F_p(t)$ and in
$\F_p[[t]]$, and the map from such fractions to series is injective; the
congruence gives $\overline F=\overline G$ in $\F_p(t)$, so
$\overline{D_F}(\alpha)=0$. Some cyclotomic factor $\Phi_d$ of $D_F$ has
$\overline{\Phi_d}(\alpha)=0$, and by Lemma~\ref{lem:cyclotomic} the prime-to-$p$
part of $d$ is $M$: $d=Mp^r$. As $N_F,D_F$ are coprime over $\Q$, $F$ has a
pole at a primitive $d$-th root of unity, and by rationality over $\Q$ at all
of them. A denominator $\prod_i(1-t^{a_i})^{b_i}$ must vanish there, so
$d\mid a_i$ for some $i$.
\end{proof}

\begin{remark}
The hypothesis on the poles of order $Mp^r$ cannot be dropped: $G(t)=1+pt^{pM}/
(1-t^{pM})$ has integral nonnegative coefficients and $\lim_{t\to\zeta}
(1-t^M)G(t)=1$, but $\overline G=1$. Nor can the conclusion be improved to
$r=0$: for $p=3$, $M=2^a\ge4$ and $z=t^{M/2}$, the series
$F=(1+z)/(1-z+z^2)+2t/(1-t)^2$ and $G=1/(1+z)+2t/(1-t)^2$ have nonnegative
integral coefficients, $F-G=3z/(1+z^3)$, the residue of $G$ at the primitive
$M$-th roots is $2$, a $3$-adic unit, but $F$ has its poles there only at order
$3M$, since $1-z+z^2=\Phi_{3M}(t)$.
\end{remark}

